\item \textbf{Cuestión de plomería.}
Está trabajando para un plomero que está colocando secciones muy largas de tubería de cobre para un gran proyecto de construcción. Pasa mucho tiempo midiendo las longitudes de las secciones con una cinta métrica. Sugiere una forma más rápida de medir la longitud. Usted sabe que la velocidad de una onda de compresión unidimensional que viaja a lo largo de una tubería de cobre es de 3.56 km/s. Usted sugiere que un trabajador dé un golpe fuerte con un martillo en un extremo de la tubería. Usando una aplicación de osciloscopio en su teléfono inteligente, medirá el intervalo de tiempo $\Delta t$ entre la llegada de las dos ondas de sonido debido al golpe: una a través del aire de 20.0 °C y la otra a través de la tubería.
\begin{enumerate}
    \item Para medir la longitud, deduzca una ecuación que relacione la longitud $L$ de la tubería numéricamente con el intervalo de tiempo $\Delta t$.
    \item Mida un intervalo de tiempo de $\Delta t = 127\text{ ms}$ entre las llegadas de los impulsos y, a partir de este valor, determine la longitud de la tubería.
    \item La aplicación de su teléfono inteligente afirma una precisión de 1.0\% en intervalos de tiempo de medición. Entonces calcule por cuántos centímetros puede haber un error en el cálculo de la longitud.
\end{enumerate}